\documentclass[../Main.tex]{subfiles}

\begin{document}
\section{Poincaré Indices}
Consider a general dynamical system:
\begin{equation*}
    \dvec{x} = \vec{f}(\vec{x}) = \begin{pmatrix}f_1(x, y) \\ f_2(x, y)\end{pmatrix}
\end{equation*}
At any given point, the angle the trajectory passing through that point makes with the horizontal is:
\begin{equation*}
    \psi(x, y) = \tan^{-1} \left(\frac{f_2(x, y)}{f_1(x, y)}\right)
\end{equation*}
Note this is not defined at fixed points.

Suppose $\Gamma$ is a simple, closed curve not passing through any fixed points. Starting at some initial point, the integral of $\psi$ along the curve (summing up all the angles along the curve) must yield $2n\pi + \phi(x_0, y_0)$. We call this $n$ the \underline{Poincaré index}:
\begin{definition}{Poincaré index}
    The \underline{Poincaré index} of a simple, closed curve $\Gamma$ not passing through any fixed points is:
    \begin{equation*}
        I(\Gamma) = \frac{1}{2\pi} \oint_\Gamma d\psi = \frac1{2\pi} \oint_\Gamma \frac{f_1 df_2 - f_2 df_1}{f_1^2 + f_2^2}
    \end{equation*}
\end{definition}
\begin{propositions}{
        Let $\Gamma$ be a simple, closed curve not passing through any fixed points.
        \label{propsPoincIndxProps}
    }
    \item $I(\Gamma)$ is unaffected by smooth, continuous deformations of $\Gamma$ as long as this deformation does not cross any fixed points. \label{propPoincIndxDeform}
    \item If $\Gamma$ does not enclose any fixed points, $I(\Gamma) = 0$. \label{propPoincIndxZero}
    \item Suppose $\Gamma = \Gamma_1 + \Gamma_2$, where $\Gamma_i$ are smooth and do not pass through fixed points. Then $I(\Gamma) = I(\Gamma_1) + I(\Gamma_2)$. \label{propPoincIndxAdd}
    \item $\Gamma$ is symmetric in $t \mapsto -t$. \label{propPoincIndxTimeReverse}
    \item If $\Gamma$ is a periodic orbit, $I(\Gamma) = 1$.
\end{propositions}
\begin{proof}
    \begin{enumerate}
        \item This is a consequence of the deformation property of integrals in $\C$. Note that the integrand (of $x + iy$) is analytic at all points other than fixed points, which means the deformation fails when crossing them.
        \item A consequence of \thmref{propPoincIndxDeform}. Deform $\Gamma$ to an infinitesimal contour on which the field is uniform, so there is no change in the index.
        \item This is simply due to the additive property of line integrals.
        \item We can see this in the definition:
            \begin{equation*}
                I(\Gamma) \mapsto \frac1{2\pi} \oint_\Gamma \frac{(-f_1) (-df_2) - (-f_2) (-df_1)}{f_1^2 + f_2^2} = \frac1{2\pi} \oint_\Gamma \frac{f_1 df_2 - f_2 df_1}{f_1^2 + f_2^2} = I(\Gamma).
            \end{equation*}
        \item On a periodic orbit $\psi$ is simply a tangent to the orbit. Evaluating the integral then gives: %TODO
    \end{enumerate}
\end{proof}
Based on these definitions, it is sensible to define:
\begin{definition}{Poincaré index of a fixed point}
    The \underline{Poincaré index} of a fixed point $\vec{x_0}$ is the Poincaré index of any curve that encircles only $\vec{x_0}$.
\end{definition}
\end{document}